\scanpage{47}
% archive sheet number: 46

\mspage{10}
\deleted{(Idem pour les $D$-équivalences \deleted{fortes}\inserted{``complètes''}, i.e. des coeff. \uline{loc} constants $\xi$ au lieu de constants. Mais il faut supposer que l'on sache déjà que $X\to\bar X$ est $D$-équi. \unsure{2-complète|2-complètes},}

Pour avoir l'implication inverse

\[
i\in W_{\underline D}\Longrightarrow i_0\in W_{\underline D},
\]

il faudrait savoir que le diagramme $(**)$, plus \deleted{\unsure{point|précis}} précisément l'é\deleted{\unsure{lément}} \deleted{diagramme}\inserted{carré} substantiel $\gamma$ qui lui donne naissance, est non seulement \uline{car}tésien, mais aussi \uline{cocartésien}. ``Moralement'', qu'il soit cartésien me semble être la traduction cohomologique du fait que le \deleted{diagramme}\inserted{carré}

\diagram{\begin{tikzcd}[ampersand replacement=\&]
\& U\arrow[dl] \&\\
U\cap U' \& \& X\arrow[ul]\arrow[dl]\\
\& U'\arrow[ul] \&
\end{tikzcd}}

\unsure{dans|ds} Cat \deleted{\unsure{est}} est \uline{cocartésien} dans un sens suffisamment fort. Mais il est également cartésien, de façon triviale, et on peut s'attendre (sous certaines conditions\ldots) que les diagrammes ``transformés'' par $H^\bullet_{\underline D}(-,\xi)$ sont cocartésiens\ldots

J'ai envie de regarder d'emblée $X$ comme sur $\underline Q=\Delta^1\times\Delta^1$, plutôt que sur \unsure{Nommons|Nommant} $\underline\Psi=\underline Q-\{e\}$ de $\underline Q$, alors

\[
\underline Y_{\underline Q}(X)|\underline\Psi\simeq Y_{\underline\Psi}(X)\ \text{déjà examiné.}
\]
